\begin{proposition}
\label{proposition-unique-base-field}
Let $X$ be a scheme. Let $a : X \to \Spec(k_1)$ and
$b : X \to \Spec(k_2)$ be morphisms from $X$ to spectra of fields.
Assume $a, b$ are locally of finite type, and
$X$ is reduced, and connected. Then we have
$k_1' = k_2'$, where $k_i' \subset \Gamma(X, \mathcal{O}_X)$ is
the integral closure of $k_i$ in $\Gamma(X, \mathcal{O}_X)$.
\end{proposition}

\begin{proof}
First, assume the lemma holds in case $X$ is quasi-compact (we will
do the quasi-compact case below).
As $X$ is locally of finite type over a field, it is locally Noetherian, see
Morphisms, Lemma \ref{morphisms-lemma-finite-type-noetherian}.
In particular this means that it is locally connected,
connected components of open subsets are open, and
intersections of quasi-compact opens are quasi-compact, see
Properties, Lemma \ref{properties-lemma-Noetherian-topology},
Topology, Lemma \ref{topology-lemma-locally-connected},
Topology, Section \ref{topology-section-noetherian}, and
Topology, Lemma \ref{topology-lemma-constructible-Noetherian-space}.
Pick an open covering $X = \bigcup_{i \in I} U_i$
such that each $U_i$ is quasi-compact and connected.
For each $i$ let $K_i \subset \mathcal{O}_X(U_i)$ be the integral
closure of $k_1$ and of $k_2$.
For each pair $i, j \in I$ we decompose
$$
U_i \cap U_j = \coprod U_{i, j, l}
$$
into its finitely many connected components. Write
$K_{i, j, l} \subset \mathcal{O}(U_{i, j, l})$
for the integral closure of $k_1$ and of $k_2$. By
Lemma \ref{lemma-integral-closure-ground-field}
the rings $K_i$ and $K_{i, j, l}$ are fields.
Now we claim that $k_1'$ and $k_2'$ both equal the kernel of the map
$$
\prod K_i \longrightarrow \prod K_{i, j, l}, \quad
(x_i)_i \longmapsto x_i|_{U_{i, j, l}} - x_j|_{U_{i, j, l}}
$$
which proves what we want.
Namely, it is clear that $k_1'$ is contained in this kernel.
On the other hand, suppose that $(x_i)_i$ is in the kernel.
By the sheaf condition $(x_i)_i$ corresponds to $f \in \mathcal{O}(X)$.
Pick some $i_0 \in I$ and let $P(T) \in k_1[T]$ be a monic polynomial
with $P(x_{i_0}) = 0$. Then we claim that $P(f) = 0$ which proves
that $f \in k_1$. To prove this we have to show that $P(x_i) = 0$
for all $i$. Pick $i \in I$. As $X$ is connected there exists a
sequence $i_0, i_1, \ldots, i_n = i \in I$ such that
$U_{i_t} \cap U_{i_{t + 1}} \not = \emptyset$. Now this means that
for each $t$ there exists an $l_t$ such that $x_{i_t}$ and
$x_{i_{t + 1}}$ map to the same element of the field $K_{i_t, i_{t + 1}, l_t}$.
Hence if $P(x_{i_t}) = 0$, then $P(x_{i_{t + 1}}) = 0$. By
induction, starting with $P(x_{i_0}) = 0$ we deduce that
$P(x_i) = 0$ as desired.

\medskip\noindent
To finish the proof of the lemma we prove the lemma under the
additional hypothesis that $X$ is quasi-compact. By
Lemma \ref{lemma-integral-closure-ground-field}
after replacing $k_i$ by $k_i'$
we may assume that $k_i$ is integrally closed in $\Gamma(X, \mathcal{O}_X)$.
This implies that $\mathcal{O}(X)^*/k_i^*$ is a finitely generated
abelian group, see
Proposition \ref{proposition-units-general}.
Let $k_{12} = k_1 \cap k_2$ as a subring of $\mathcal{O}(X)$.
Note that $k_{12}$ is a field. Since
$$
k_1^*/k_{12}^* \longrightarrow \mathcal{O}(X)^*/k_2^*
$$
we see that $k_1^*/k_{12}^*$ is a finitely generated abelian group as well.
Hence there exist $\alpha_1, \ldots, \alpha_n \in k_1^*$ such that
every element $\lambda \in k_1$ has the form
$$
\lambda = c \alpha_1^{e_1} \ldots \alpha_n^{e_n}
$$
for some $e_i \in \mathbf{Z}$ and $c \in k_{12}$.
In particular, the ring map
$$
k_{12}[x_1, \ldots, x_n, \frac{1}{x_1 \ldots x_n}] \longrightarrow k_1, \quad
x_i \longmapsto \alpha_i
$$
is surjective. By the Hilbert Nullstellensatz,
Algebra, Theorem \ref{algebra-theorem-nullstellensatz}
we conclude that $k_1$ is a finite extension of $k_{12}$.
In the same way we conclude that $k_2$ is a finite extension of $k_{12}$.
In particular both $k_1$ and $k_2$ are contained in the integral closure
$k_{12}'$ of $k_{12}$ in $\Gamma(X, \mathcal{O}_X)$. But since $k_{12}'$
is a field by
Lemma \ref{lemma-integral-closure-ground-field}
and since we chose $k_i$ to be integrally closed in $\Gamma(X, \mathcal{O}_X)$
we conclude that $k_1 = k_{12} = k_2$ as desired.
\end{proof}